% Biographical notes. Short lives of Gene and Joan, plus a service note.
% Drawn from the family research document and the letters themselves.
% Facts here should be checked against primary records before any public
% printing; the marriage date and exact discharge date in particular.

\cleardoublepage
\thispagestyle{laempty}
\vspace*{0.12\textheight}
\begin{center}
    \eyebrow{Lives}\par
    \vspace{0.4em}
    \brassrule[36pt]\par
    \vspace{1.6em}
    {\Huge\displayfont\itshape Biographical Notes\par}
\end{center}
\vfill
\cleardoublepage

\thispagestyle{laplain}
\vspace*{0.06\textheight}
\begin{center}
    {\Large\displayfont\itshape Raymond Eugene Lankford\par}
    {\monofont\small\addfontfeature{LetterSpace=18}\color{walnutmute}1921\,--\,2008}
\end{center}
\vspace{1.4em}

\noindent
Raymond Eugene Lankford\,---\,Gene\,---\,was born on November~21, 1921, in
Stanford, Lincoln County, Kentucky, a son of William Porter Lankford and
Marguerite Lynn Lankford. He grew up in a large family; an older brother,
Carroll, had already served in the Navy when Gene, at eighteen, enlisted in
the spring of 1940 and left Kentucky for the U.S. Naval Training Station at
Great Lakes, Illinois. The first letters in this book are written from there.

\par
After training he was assigned to the heavy cruiser U.S.S. New Orleans
(CA\,--\,32), serving in her F~Division. The ship was moored at Pearl Harbor
on the morning of December~7, 1941. Over the next year she fought through the
Pacific\,---\,the Coral Sea, Midway, and the long battle for Guadalcanal. On
the night of November~30, 1942, off Tassafaronga in the Solomon Islands, a
Japanese torpedo struck the New Orleans and tore away roughly a hundred and
fifty feet of her bow. One hundred and eighty-three of her crew were killed.
Most of Joan's letters to Gene, kept aboard ship, went down with the bow; that
is why this collection runs in one direction.

\par
Gene was honorably discharged in 1943 and, on April~20 of that year, married
Joan Northcutt, the girl he had courted across four years of letters. He came
home to Kentucky, made his life in Somerset, and for many years owned and ran
Crest Electronics there. He was a member of High Street Baptist Church. He and
Joan were married for forty-nine years, until her death in 1992; he later
married Betty Pitman Lankford. Gene died at home in Somerset on May~23, 2008,
at the age of eighty-six, and was buried with military honors at Lakeside
Memorial Gardens. He was survived by his children Raymond Eugene Lankford~Jr.,
Robert Lee Lankford, and Elizabeth ``Libby'' Lankford Morris, and was
predeceased by his son William Joe Lankford and by two infant daughters,
Debra Ann and Janie Wood.

\vspace{2em}
\begin{center}
    {\Large\displayfont\itshape Joan Northcutt Lankford\par}
    {\monofont\small\addfontfeature{LetterSpace=18}\color{walnutmute}1925\,--\,1992}
\end{center}
\vspace{1.4em}

\noindent
Joan Northcutt was born on May~17, 1925, and grew up on Route~1 outside
Stanford, Kentucky. She was still a schoolgirl when Gene left for the Navy in
1940\,---\,fourteen years old\,---\,and his letters follow her through her
exams, her grades, and the small news of a Kentucky town at the edge of a war.
For the length of the war she was at home in Stanford while Gene was at sea.

\par
She kept every letter Gene wrote her, for the rest of her life. In later years
she would not let her own daughter read them; they were, she said, too
personal. She married Gene on April~20, 1943, and they were together
forty-nine years. Joan died on November~19, 1992, and is buried at Lakeside
Memorial Gardens in Somerset. The letters were found, years after her death, in
a box in a storage building behind the family home\,---\,which is how they came
to be gathered here, and given to the ship's namesake city.

\vspace{2em}
\begin{center}
    \eyebrow{A note on his service record}\par
    \vspace{0.3em}
    \brassrule[22pt]\par
\end{center}
\vspace{1em}

\noindent
Gene's exact rate and rating at discharge are not recorded in any public
source consulted for this book. They can be obtained from his official
military personnel file, held at the National Personnel Records Center in
St.~Louis and requestable by next of kin on Standard Form~180. The ship's own
record is well documented: the U.S.S. New Orleans earned seventeen battle
stars in the Pacific, and her severed bow was found on the floor of Ironbottom
Sound in 2025.
\vfill
